\chapter{The summation cover: what is fixed for every curve}\label{ch:cover}

Every index calculus on an elliptic curve asks one question many times:
does a target $R$ decompose as a sum of $m-1$ factor-base points? Semaev's
summation polynomial answers it algebraically, and every design lever of
the programme acted on this question. This chapter collects three facts
about it that hold for \emph{every} curve. They closed, without an experiment, the levers
that rely on curve-dependent monodromy, on fibre censuses, or on raising
mean yield at matched base size. They also mark what stays free
(\S\ref{sec:levers}).

The statements and proofs are those of
\art{crypto-autoresearcher:papers/semaev-conservation-specialization/}
(15 pages, anonymized draft, verification script
\texttt{verify.py}, seed 20260807, which replays byte for byte), and of the
internal finding \texttt{KN-FIND-bdcd60}. An integrity audit
(\art{crypto-autoresearcher:research/sheriff/SHF-20261023-7b9227/}) upheld
the mathematics within its scope. It blindly re-derived the group orders
$2,24,192,1920,23040$ and checked $2.1$ million admissible fibres with zero
violations. It also rated as overclaimed a statement that three
\emph{independent} derivations agree: they came from one model family.
\textbf{Novelty is not adjudicated} (\S\ref{sec:cover-prior}).

\section{Setting}

Let $K$ be a field, $E/K$ an elliptic curve with inversion $\iota=[-1]$,
and $\pi:E\to E/\langle\iota\rangle\cong\PP^1$ the quotient (the
$x$-coordinate in odd characteristic). The $m$-th summation polynomial
$\Sem m\in K[X_1,\dots,X_m]$ is the reduced defining polynomial of the
image of $\{(Q_1,\dots,Q_m):\sum Q_i=O\}$ under $\pi^m$. Equivalently,
$\Sem m(\xi_1,\ldots,\xi_m)=0$ if and only if there are points $Q_i$ with
$\pi(Q_i)=\xi_i$ summing to $O$~\cite{semaev2004}. For $m\ge3$ it is
symmetric, absolutely irreducible, and of degree $2^{m-2}$ in each
variable. It satisfies the resultant recursion
$\Sem m=\operatorname{Res}_X(\Sem{m-1}(\ldots,X),\Sem3(X_{m-1},X_m,X))$.

Fix $m\ge3$, put $n=m-1$, and let $t_1,\dots,t_n$ be independent
transcendentals. Write $L=K(t_1,\dots,t_n)$ for the function field of
$(\PP^1)^n$ and $M=K(E^n)$ for that of $E^n$. The group
$\Gamma=\langle\iota\rangle^n\cong(\ZZ/2)^n$ acts on $E^n$
coordinatewise; $\Delta\subset\Gamma$ is its diagonal. $M/L$ is Galois with
group $\Gamma$. With $P_i$ the generic points ($\pi(P_i)=t_i$), the roots of
$\Sem m(t_1,\ldots,t_n,T)$ in $T$ are the abscissae of the signed sums
$\theta_\varepsilon=\pi(\varepsilon_1P_1+\cdots+\varepsilon_nP_n)$,
$\varepsilon\in\{\pm1\}^n$, with $\theta_\varepsilon=\theta_{-\varepsilon}$.

\section{The Galois group}

\begin{theorem}[Monodromy of the summation cover]\label{thm:cover}
For every field $K$, every elliptic curve $E/K$ and every $m\ge3$:
\begin{enumerate}
\item $\Sem m(t_1,\dots,t_n,T)$ is separable of degree $2^{m-2}$ over $L$,
  with root set $\{\theta_\varepsilon\}$;
\item its splitting field is $N=M^\Delta=K(E^n/\Delta)$;
\item $N/L$ is Galois with group $\Gamma/\Delta\cong(\ZZ/2)^{m-2}$, acting
  simply transitively on the roots;
\item $K$ is algebraically closed in $N$, so arithmetic and geometric
  monodromy coincide, $\mathrm{Mon}^{\mathrm{arith}}=\mathrm{Mon}^{\mathrm{geom}}\cong(\ZZ/2)^{m-2}$.
\end{enumerate}
None of the four statements depends on $E$, on $K$ or on its
characteristic.
\end{theorem}

\begin{proof}[Proof outline]
A root-identification lemma shows that the roots are among the
$\theta_\varepsilon$. A generic-separability lemma shows that the
$2^{m-2}$ classes $\{\pm\varepsilon\}$ give distinct values. $\Gamma$
permutes the roots by $\sigma\cdot\theta_\varepsilon=\theta_{\sigma\varepsilon}$,
and the kernel of this action is exactly $\Delta$. So $\Gamma/\Delta$ acts
faithfully and regularly on $2^{m-2}$ roots. The stabiliser of one root is
$\Delta$, which gives $N=M^\Delta$. Finally, $E^n/\Delta$ is geometrically
integral, so $K$ is algebraically closed in its function field.
\end{proof}

For a symmetric (unordered) factor base the relevant group is the wreath
extension $(\ZZ/2)^{m-2}\rtimes S_{m-1}$, of order $2^{m-2}(m-1)!$ for
$m\ge4$. It is primitive if and only if $m-1$ is odd
(\texttt{KN-FIND-bdcd60}). Statement (4) is the operative one: a cover
whose arithmetic monodromy exceeds its geometric monodromy can split
differently over different base fields, and hence on different curves.
Here no such room exists.

\section{An exact factorisation law over a finite field}

Over $K=\FF_q$, define for $a\in\PP^1(\FF_q)$ the \emph{splitting
invariant}:
\begin{itemize}
\item $c(a)=0$ if both points above $a$ are rational;
\item $c(a)=1$ if neither is;
\item $c(a)=\ast$ at a branch point.
\end{itemize}
For odd $q$ and $E:y^2=f(x)$ this is the quadratic character of $f(a)$.

\begin{theorem}[Factorisation law]\label{thm:factor}
Let $a=(a_1,\ldots,a_n)\in\FF_q^n$ lie outside the degenerate locus, which
has $O_m(q^{n-1})$ points. Then $\Sem m(a_1,\ldots,a_n,T)$ is squarefree of
degree $2^{m-2}$, and:
\begin{enumerate}
\item if $c(a_1)=\cdots=c(a_n)$, it splits into $2^{m-2}$ distinct
  linear factors, with roots $x(\pm P_1\pm\cdots\pm P_n)$;
\item otherwise it is a product of $2^{m-3}$ distinct irreducible
  quadratics and has no root in $\FF_q$.
\end{enumerate}
No irreducible factor of degree $\ge3$ ever occurs.
\end{theorem}

\begin{proof}
On the étale locus, the factorisation type is the cycle type of
Frobenius in $\Gamma/\Delta$ (Dedekind's theorem for specialisations of a
Galois cover~\cite{friedjarden2008}). The Frobenius element is
$((-1)^{c(a_i)})_i$. It is trivial in $\Gamma/\Delta$ exactly when every
$c(a_i)$ agrees. Otherwise it is a non-trivial element of order 2 of a
regular group, and so has cycle type $2^{2^{m-3}}$.
\end{proof}

The density of totally split fibres is $2^{-(m-2)}+O(m^2/(2^{m-1}q))$. The
first-order Hasse terms $\pm t/q$ cancel, so the statistic carries no
first-order information about the curve.

\section{Specialisation to the factor-base locus}

\begin{theorem}[Complete splitting on the factor-base locus]\label{thm:fbsplit}
If each $a_i=x(P_i)$ with $P_i\in E(\FF_q)$ and $a$ is non-degenerate, then
$\Sem m(a_1,\ldots,a_n,T)$ splits into $2^{m-2}$ distinct linear factors
over $\FF_q$, with roots $\{x(\varepsilon_1P_1+\cdots+\varepsilon_nP_n)\}$.
\end{theorem}

This is Theorem~\ref{thm:factor}(1) with every $c(a_i)=0$. It is the
statement behind the root-finding oracle of Chapter~\ref{sec:rootfinding}:
given $m-1$ base points, the admissible last coordinates are the signed
sums, computable in $O(2^{m-2})$ group operations. It is \emph{not} an
algorithm for relation search. Relation search fixes the target and asks
for the base points. The theorem says nothing about how many tuples must
be tried, nor about the algebraic difficulty of the system posed in that
direction. Its value is negative:

\begin{corollary}[Census obstruction]\label{cor:census}
Every statistic of the form ``factorisation type, cycle type or Frobenius
class of the fibre of $\Sem m$'' is constant on the factor-base locus, for
every curve and every $m$. A census of generic fibres measures a quantity
whose value where relation search operates is known in advance and is the
same for all curves.
\end{corollary}

The programme ran such censuses before this corollary was proved. The
census-as-relation-rate probe was closed by ledger decision
\texttt{DEC-20261017-435b01} (\texttt{reject\_scoped}, replicated).

\section{Conservation of decomposition yield}

\begin{proposition}[Conservation]\label{prop:conservation}
Let $G$ be a finite abelian group of order $N$ and
$\mathcal D\subseteq G\setminus\{0\}$ a set of size $B$. The number of
size-$m$ multisets from $\mathcal D$, summed over all targets, is
$\binom{B+m-1}{m}$. The mean number of decompositions per target is
therefore $\binom{B+m-1}{m}/N$, for every $\mathcal D$ of size $B$.
\end{proposition}

The proof is a double count, and no novelty is claimed for it. Its
consequences shaped the programme:
\begin{itemize}
\item at matched base size, no factor-base geometry can raise the mean
  yield;
\item only the distribution of decompositions is free, and coverage
  (the fraction of targets with at least one decomposition) is at most
  $\min(1,\mu)$;
\item additive structure that repeats sums \emph{lowers} coverage;
\item typing the summands, each from its own sub-base, is a fixed penalty:
  $\prod B_i\le(B/m)^m<\binom{B+m-1}{m}$.
\end{itemize}
Gaudry and Diem's $1/m!$ decomposition probability over extension fields
\emph{is} this mean. The extension-field advantage lives in the solve, not
in the yield.

\section{Which levers are fixed and which are free}\label{sec:levers}

Together the results above divide the design space as follows:
\begin{itemize}
\item \emph{Fixed for every curve:}
  \begin{itemize}
  \item the monodromy group of the summation cover (Theorem~\ref{thm:cover});
  \item the fibre statistics on the factor-base locus
    (Corollary~\ref{cor:census});
  \item the mean yield at fixed base size (Proposition~\ref{prop:conservation}).
  \end{itemize}
\item \emph{Free:}
  \begin{itemize}
  \item the coverage distribution (bounded by the mean);
  \item the cost of deciding membership (the oracle, Part~\ref{part:algorithms});
  \item the algebraic difficulty of the decomposition system when it is
    posed as a polynomial system (Chapter~\ref{ch:falldegree});
  \item the linear algebra;
  \item where the problem lives (transfers, Chapter~\ref{sec:transfer}).
  \end{itemize}
\end{itemize}
Every measured improvement in Part~\ref{part:algorithms} acts on a free
item. Chapter~\ref{sec:transfer} shows that transfers by isogeny conserve
the counts as well.

\section{Prior work and the limits of the novelty assessment}\label{sec:cover-prior}

Absolute irreducibility of $\Sem m$ is classical~\cite{semaev2004,gaudry2009}.
Corollary~\ref{cor:census}'s ingredient, that the roots are the signed sums,
is the defining property of $\Sem m$. The action of 2-torsion and of sign
changes on summation systems is used by Faugère, Gaudry, Huot and
Renault~\cite{fghr2014} and by Galbraith and
Gebregiyorgis~\cite{galbraith2014}. Theorem~\ref{thm:cover}'s group may be
folklore: it is the deck group of $E^n\to(E/\pm)^n$ modulo the diagonal.
Theorem~\ref{thm:factor} is in essence a 2-descent statement. The draft
could not reach ePrint during its literature search, and says so.
Corollary~\ref{cor:census} and the exact dichotomy of
Theorem~\ref{thm:factor} are the parts most likely to be new. Two items
must be settled before any of this is submitted:
\begin{itemize}
\item an expert check of the separability lemma in characteristic 2,
  where $E[2]$ is non-reduced;
\item a full literature search.
\end{itemize}
